\section*{الفصل \ref{ch:g12:exp} --- الدالة الأسية واللوغاريتم}

\begin{solution}{exo:g12:exp:1}
$\dfrac{\eu^{3x}\,\eu^{-x+1}}{\eu^{x}} = \eu^{3x - x + 1 - x} = \eu^{x+1}$.

$\ln\!\left(\dfrac{\eu^2\sqrt{\eu}}{\eu^{-3}}\right)
= 2 + \tfrac12 + 3 = \tfrac{11}{2}$.
\end{solution}

\begin{solution}{exo:g12:exp:2}
\emph{(أ)} ضع $X = \eu^x > 0$: فإن $X^2 - 3X + 2 = 0$ يعطي $X = 1$ أو $X = 2$،
وكلاهما موجب، إذن $x = 0$ أو $x = \ln 2$.

\emph{(ب)} \omterm{def:g11:func:function}{مجموعة التعريف}: $x - 1 > 0$ و $x + 2 > 0$، إذن $x > 1$. وتصير \omterm{def:g10:algebra:equation}{المعادلة}
$\ln\bigl((x-1)(x+2)\bigr) = \ln 4$، ومنه $(x-1)(x+2) = 4$،
\emph{أي} $x^2 + x - 6 = 0$، إذن $x = 2$ أو $x = -3$. ولا يقع في \omterm{def:g11:func:function}{مجموعة التعريف}
إلا $x = 2$: فيكون $S = \{2\}$.
\end{solution}

\begin{solution}{exo:g12:exp:3}
اقسم على $\eu^x$:
$\dfrac{1 - x^2\eu^{-x}}{1 + \eu^{-x}} \to \dfrac{1-0}{1+0} = 1$، باستعمال
$x^2 \eu^{-x} \to 0$ (\cref{thm:g12:exp:growth}).

$\dfrac{\ln x}{\sqrt x} \to 0$ حسب \cref{prop:g12:exp:lnanalytic} (الحالة
$n = 2$).

$x^2 \ln x = x \cdot (x \ln x) \to 0 \times 0 = 0$.

$x - \ln x = x\left(1 - \frac{\ln x}{x}\right) \to +\infty$ لأن
$\frac{\ln x}{x} \to 0$.
\end{solution}

\begin{solution}{exo:g12:exp:4}
$f'(x) = \eu^{-x} - x\eu^{-x} = (1 - x)\eu^{-x}$، وإشارتها إشارة $1 - x$:
فتتزايد $f$ على $\intoc{-\infty}{1}$، وتتناقص على $\intco{1}{+\infty}$،
مع \omterm{def:g10:functions:extrema}{قيمة عظمى} شاملة $f(1) = \eu^{-1}$.

والنهايات: لما $x \to +\infty$، $f(x) = \frac{x}{\eu^x} \to 0$
(\cref{thm:g12:exp:growth})؛ ولما $x \to -\infty$، يكون $x \to -\infty$ و
$\eu^{-x} \to +\infty$، إذن $f(x) \to -\infty$.

$f''(x) = -\eu^{-x} - (1-x)\eu^{-x} = (x - 2)\eu^{-x}$، وهي تغيّر إشارتها عند
$x = 2$: \omterm{def:g12:deriv:inflection}{فنقطة الانعطاف} عند $\bigl(2,\, 2\eu^{-2}\bigr)$.
\end{solution}

\begin{solution}{exo:g12:exp:5}
لتكن $g(x) = x - \ln(1+x)$ على $\intoo{-1}{+\infty}$. عندئذٍ
$g'(x) = 1 - \frac{1}{1+x} = \frac{x}{1+x}$، وهي سالبة على $\intoo{-1}{0}$ و
موجبة على $\intoo{0}{+\infty}$: فتبلغ $g$ قيمتها الصغرى $g(0) = 0$، إذن
$g \geq 0$ مع المساواة عند $0$ فقط. ومنه $\ln(1+x) \leq x$.

طبّق ذلك مع $x = \frac1n$:
$n\ln\left(1 + \frac1n\right) \leq 1$، إذن
$\left(1+\frac1n\right)^n = \eu^{\,n\ln(1+1/n)} \leq \eu$.
وطبّقه مع $x = -\frac{1}{n+1} > -1$:
$\ln\left(1 - \frac{1}{n+1}\right) \leq -\frac{1}{n+1}$، إذن
$-(n+1)\ln\left(1 - \frac{1}{n+1}\right) \geq 1$ و
$\left(1 - \frac{1}{n+1}\right)^{-(n+1)} \geq \eu$.
\end{solution}

\begin{solution}{exo:g12:exp:6}
\emph{1.} كل سنة تضرب رأس المال في $1.03$: $C_n = C_0 (1.03)^n$.

\emph{2.} $C_n \geq 2C_0 \iff (1.03)^n \geq 2 \iff n \ln 1.03 \geq \ln 2
\iff n \geq \frac{\ln 2}{\ln 1.03} \approx 23.45$: فيتضاعف رأس المال بعد
$24$ سنة.

\emph{3.} $\frac{70}{3} \approx 23.3$، وهو قريب من القيمة المضبوطة
$\frac{\ln 2}{\ln 1.03}$. وتنجح القاعدة لأن
$\ln(1 + r) \approx r$ من أجل $r$ صغير، إذن
$\frac{\ln 2}{\ln(1+r)} \approx \frac{0.693}{r} \approx \frac{70}{100r}$.
\end{solution}

\begin{solution}{exo:g12:exp:7}
$\eu^{2x} > \eu^{x+1} \iff 2x > x + 1 \iff x > 1$ (\omterm{thm:g12:exp:existence}{فالدالة الأسية}
\omterm{def:g12:seq:monotonic}{متزايدة} تمامًا): $S = \intoo{1}{+\infty}$.

\omterm{def:g10:functions:function}{ومجموعة تعريف} المتراجحة الثانية: $x^2 - 1 > 0$ و $x + 5 > 0$، \emph{أي}
$x \in \intoo{-5}{-1} \cup \intoo{1}{+\infty}$. وعلى هذه المجموعة، وبما أن $\ln$
\omterm{def:g12:seq:monotonic}{متزايدة} تمامًا،
\[
\ln(x^2-1) \leq \ln(x+5) \iff x^2 - 1 \leq x + 5 \iff x^2 - x - 6 \leq 0
\iff x \in \intcc{-2}{3}.
\]
وبالتقاطع مع \omterm{def:g11:func:function}{مجموعة التعريف}: $S = \intco{-2}{-1} \cup \intoc{1}{3}$.
\end{solution}

\begin{solution}{exo:g12:exp:8}
\emph{1.} $f'(x) = \dfrac{\eu^x x - \eu^x}{x^2} = \dfrac{(x-1)\eu^x}{x^2}$،
وهي سالبة على $\intoo{0}{1}$ وموجبة على $\intoo{1}{+\infty}$: \omterm{def:g10:functions:extrema}{فالقيمة الصغرى}
$f(1) = \eu$. والنهايات: $f \to +\infty$ عند $0^+$ (البسط $\to 1$،
والمقام $\to 0^+$) وعند $+\infty$ (\cref{thm:g12:exp:growth}).

\emph{2.} من أجل $x > 0$، $\eu^x = kx \iff f(x) = k$. ومن \omterm{def:g10:functions:table}{جدول التغيرات}
($+\infty \searrow \eu \nearrow +\infty$، \omterm{def:g12:limcont:continuity}{متصلة} \omterm{def:g12:seq:monotonic}{ورتيبة} تمامًا
على كل جهة): لا حل من أجل $k < \eu$؛ وحل واحد بالضبط ($x = 1$) من أجل
$k = \eu$؛ وحلان بالضبط من أجل $k > \eu$ (أحدهما في $\intoo{0}{1}$ والآخر في
$\intoo{1}{+\infty}$، بمبرهنة التقابل على كل \omterm{def:g10:numbers:interval}{فترة}).
\end{solution}

\begin{solution}{exo:g12:exp:9}
\emph{1.} مباشرة من المتراجحة الأولى في \cref{exo:g12:exp:5}:
$u_n = \eu^{\,n\ln(1+1/n)} \leq \eu^1 = \eu$.

\emph{2.} اكتب
\[
\ln u_n = n \ln\!\left(1 + \frac1n\right)
= \frac{\ln\!\left(1 + \frac1n\right) - \ln 1}{\frac1n},
\]
وهو معدل تغير \omterm{def:g11:func:function}{للدالة} $\ln$ عند النقطة $1$ بزيادة
$h = \frac1n \to 0$. وبما أن $\ln$ قابلة للاشتقاق عند $1$ ومشتقتها
$1$، فإن $\ln u_n \to 1$. وباتصال $\exp$
(\cref{prop:g12:limcont:seqcont})، $u_n = \eu^{\ln u_n} \to \eu^1 = \eu$.
\end{solution}

\begin{solution}{pb:g12:exp:1}
\textbf{1.} $x = \frac{\ln 5}{2}$. \quad
$3x - 1 = \eu^2$: $x = \frac{\eu^2 + 1}{3}$. \quad ومع
$u = \eu^x$: $u^2 - 3u + 2 = (u - 1)(u - 2) = 0$: أي $x = 0$ أو
$x = \ln 2$.

\textbf{2.} $\frac{3\ln 2}{\ln 2} = 3$؛ \quad
$3 + \frac12 = \frac72$؛ \quad $\frac52$.

\textbf{3.} $f'(x) = \eu^x - 1$: سالبة قبل $0$ وموجبة
بعدها: \omterm{def:g10:functions:extrema}{فالقيمة الصغرى} $f(0) = 0$، إذن $f \geq 0$ في كل مكان:
$\eu^x \geq 1 + x$. وبالتعويض $x = \ln(1 + u)$:
$1 + u \geq 1 + \ln(1 + u)$، أي $\ln(1 + u) \leq u$.

\textbf{4.} $x^2\eu^{-x} = \frac{x^2}{\eu^x} \to 0$ و
$\frac{\ln x}{x} \to 0$: فالدوال الأسية تسحق القوى، والقوى
تسحق اللوغاريتمات. و $\frac{\eu^x - 1}{x} =
\frac{\eu^x - \eu^0}{x - 0} \to \exp'(0) = 1$.

\textbf{5.} $f'(x) = \ln x + 1$: تنعدم عند $\eu^{-1}$، وهي سالبة
قبله وموجبة بعده: \omterm{def:g10:functions:extrema}{فالقيمة الصغرى}
$f\!\left(\frac1\eu\right) = -\frac1\eu$؛ و
$x \ln x \to 0$ عند $0^+$: \omterm{def:g10:functions:graph}{فالمنحنى} يغادر المبدأ، وينخفض إلى
$-\frac1\eu$، ثم يتسلق مبتعدًا.

\textbf{6.} $n = \frac{\ln 2}{\ln 1.03} \approx 23.4$ سنة.

\textbf{7.} المضبوط: $23.4$ و $11.9$ و $8.0$ سنة؛ وقاعدة 72:
$24$ و $12$ و $8$. وبما أن $\ln(1+r) \approx r$، فإن
$\frac{\ln 2}{\ln(1+r)} \approx \frac{0.693}{r}$، أي
$\frac{69.3}{\text{المعدل}\%}$؛ ويقرّب المصرفيون إلى $72$
لأنه يقبل القسمة قسمة جميلة على $2, 3, 4, 6, 8, 9, 12$ ---
فالحساب الذهني يتغلب على رقم عشري من الدقة.

\textbf{8.} $2^{t/20} = \eu^{t \ln 2 / 20}$:
أي $\lambda = \frac{\ln 2}{20}$ في الدقيقة. وبعد
$24 \times 60 = 1440$ دقيقة: $2^{72} \approx 4.7 \times
10^{21}$ بكتيريا --- أي أكثر من حبات الرمل على الأرض،
من خلية واحدة في يوم واحد. والنموذج أمين ما دام الغذاء
والمكان متوافرين: أما النمو الحقيقي فينحني إلى \omterm{def:g10:functions:graph}{منحنى} S اللوجستي
(وهو \omterm{def:g10:functions:graph}{منحنى} عالِم الأوبئة في \cref{pb:g12:deriv:1}).

\textbf{9.} عند 23:00 (أي $8$ ساعات):
$100 \times 2^{-8/5} \approx 33$ mg. وتحت $10$ mg:
يعطي $2^{-t/5} < 0.1$ أن $t > 5\,\frac{\ln 10}{\ln 2} \approx
16.6$ h: أي نحو 7:40 من صباح اليوم التالي --- فللإسبريسو عند
الثالثة ذيل طويل.

\textbf{10.} سنويًا: $2$. وشهريًا:
$\left(1 + \frac{1}{12}\right)^{12} \approx 2.613$. ويوميًا:
$\approx 2.715$. والسقف: $\eu = 2.71828\ldots$
(\cref{exo:g12:exp:9}) --- فبالتركيب المتصل، يتقارب
الجشع.

\textbf{11.} إذا نمت الحالات نموًا أسيًا، $c(t) = c_0
\eu^{\lambda t}$، فإن $\ln c(t) = \ln c_0 + \lambda t$: أي
مستقيم ميله $\lambda$ --- وهو معدل النمو. والامتداد
المستقيم على الرسم اللوغاريتمي \emph{هو} الطور
الأسي، وانحداره هو إيقاع الوباء.

\textbf{12.} $120 - 60 = 60$ dB يعني
$10\log_{10}(I_2/I_1) = 60$: فالحفل أشد من
المحادثة $10^6$ --- أي مليون مرة.

\textbf{13.} السعة: $10^{7-5} = 100$. والطاقة:
$100^{3/2} = 1000$: فدرجتان من المقياس تخفيان ثلاث رتب
قدر في الطاقة.

\textbf{14.} $10^{6.5 - 2} = 10^{4.5} \approx 32\,000$: فعصير الليمون
أشد حموضة من الحليب ثلاثين ألف مرة --- إذ يضغط مقياس
الحموضة هوّات الكيمياء في سلّم \omterm{def:g11:trigo:cossin}{جيب}.

\textbf{15.} $\log_2\!\frac{4186}{27.5} =
\frac{\ln(4186/27.5)}{\ln 2} \approx 7.25$: فالبيانو يمتد على ما يزيد
قليلًا على سبعة أوكتافات. والحواس تستجيب إلى \emph{نسب} المنبهات
--- فمضاعفة الشدة تُحسّ خطوة واحدة، مهما كان
المستوى الابتدائي --- إذن الإدراك مبني على سلّم
لوغاريتمي، وكذلك الوحدات التي اخترعناها له.

\textbf{16.} وضع عصا العدد $a$ طرفًا إلى طرف مع
موضع العدد $b$ يجمع الطولين $\ln a + \ln b$، و
التدريجة الجالسة عند ذلك الطول الكلي تُقرأ
$\eu^{\ln a + \ln b} = ab$: فالمسطرة الحاسبة تحسب الجداءات
بجمع اللوغاريتمات --- إنها $\ln(ab) = \ln a + \ln b$
(\cref{prop:g12:exp:lnalg}) محفورة في خشب البقس.

\textbf{17.} \omterm{def:g10:functions:extrema}{القيمة الصغرى} \omterm{def:g11:func:function}{للدالة} $\frac{\eu^x}{x}$ على
$\intoo{0}{+\infty}$ هي $\eu$، عند $x = 1$: فلا حل من أجل
$k < \eu$، وحل واحد بالضبط من أجل $k = \eu$، وحلان من أجل $k > \eu$.
والتحقق \omterm{def:g12:deriv:tangent}{بالمماس}: عند $a = 1$ يكون \omterm{def:g12:deriv:tangent}{المماس} \omterm{def:g10:functions:graph}{للمنحنى} $\eu^x$ هو
$y = \eu + \eu(x - 1) = \eu x$: وهو يمر بالمبدأ
--- إنه مستقيم العتبة، الذي يلامس \omterm{def:g10:functions:graph}{المنحنى} عند $(1, \eu)$.

\textbf{18.} $n \ln\!\left(1 - \frac1n\right) =
\frac{\ln(1 - 1/n)}{-1/n} \times (-1) \to -1$، إذن
$\left(1 - \frac1n\right)^n \to \eu^{-1} \approx 0.368$: أي
$0.999^{1000} \approx 0.368$ في \cref{pb:g11:binom:1}،
مفسَّرًا --- فثابت الفوز الضائع في اليانصيب هو
$\frac1\eu$.

\textbf{19.} كلاهما الشكل النهائي لعبارة ``أحداث كثيرة مستقلة،
كل منها في ذاته غير محتمل'': \omterm{def:g10:proba:distribution}{فاحتمال} ألا ينطلق
\emph{أي} من $n$ فرصة حجم كل منها $\frac1n$ يؤول إلى
$\eu^{-1}$، وقاعدة السكرتيرة تضبط نافذة رفضها
حتى يتركز النجاح عند ذلك الثابت نفسه.
$\frac1\eu \approx 0.368$.

\textbf{20.} سقف التركيب؛ \omterm{def:g11:seq:arithmetic}{والأساس} الوحيد الذي
\omterm{def:g10:reffunc:affine}{ميل} مماسه عند $0$ يساوي $1$ (ولهذا يحبه التحليل)؛
والنمو الذي يسبق كل قوة؛ والثابت الذي عنده تستقر
حالات الفوز الضائع والتوقف الأمثل؛ \omterm{def:g12:exp:ln}{واللوغاريتم} ---
الضرب وقد صار جمعًا، \omterm{def:g10:stats:quartiles}{ومدى} العالم الأوسع وقد
طُوي على مسطرة.
\end{solution}
